\section{Periodic fields and a prescribed continuum law}
\label{sec:continuum-limits}

A finite periodic step field does not satisfy the sharp global
$L^2(\R^3)$ Gaussian bound for arbitrary Schwartz tests.
The periodized norm must be retained until the volume tends
to infinity. We give the moment argument of
\cite[\S6.1--6.2]{QED3Master}; it is the step needed to use the
finite-volume product estimates for a prescribed parent.

Let $\mathsf A=1+|y|^2-\Delta$ on $\R^3$, with real
$L^2$-orthonormal Hermite two-forms $h_\kappa$ and eigenvalues
$\lambda_\kappa$. Each scalar Hermite function occurs once
for each of the three components. We use the Hilbert space
\[
 H_{-4}=\left\{F:\sum_\kappa\lambda_\kappa^{-8}
                                  |F(h_\kappa)|^2<\infty\right\}.
\]
It is continuously embedded in tempered two-form distributions,
and each Schwartz evaluation is continuous there.
The step-field realization, tightness and fixed-spacing limits
in this space are proved in \cite{ContinuumExistence}.

\subsection{The periodic seminorm}

For $R=n\eta$ and a Schwartz two-form $g$, set
$g^R(y)=\sum_{k\in\Z^3}g(y+Rk)$ and
\[
 g_\eta([\sigma])=\eta^{-3}
                \int_{Q_\eta(\sigma)}g^R_{\mu\nu}(y)\,\dd y,
 \qquad |g_\eta|_\eta^2=\eta^3\sum_{[\sigma]}|g_\eta([\sigma])|^2.
\]
Unfolding the absolutely convergent sum gives
\begin{equation}\label{eq:limit-step-pairing}
 F^\sharp(g)=\eta^{3/2}\sum_{[\sigma]}
                         g_\eta([\sigma])F^{\rm u}(\sigma).
\end{equation}
Cell Cauchy--Schwarz and Minkowski give
\begin{equation}\label{eq:limit-torus-norm}
 |g_\eta|_\eta\le\|g^R\|_{L^2(\R^3/R\Z^3)}\le p_R(g),
 \qquad
 p_R(g)=\|g\|_2+
      26R^{-2}\||y|^2g\|_{L^2(\R^3\setminus Q_R)},
\end{equation}
where $Q_R=[-R/2,R/2)^3$.
For completeness, put $a_k=\|g\|_{L^2(Q_R+Rk)}$.
The periodized norm is at most $\sum_ka_k$. On a cube with
$k\ne0$, $|y|\ge R|k|_\infty/2$, so
\[
 \sum_{k\ne0}a_k
 \le4R^{-2}\left(\sum_{k\ne0}|k|_\infty^{-4}\right)^{1/2}
                 \||y|^2g\|_{L^2(\R^3\setminus Q_R)}.
\]
The shell count $24j^2+2$ gives a sum
$4\pi^2+\pi^4/45<41.65$; its square root times four is
less than $26$. This proves \eqref{eq:limit-torus-norm}.
For $R\ge1$,
\[
 p_R(g)\le p_1(g),\qquad p_R(g)\longrightarrow\|g\|_2
                      \quad(R\longrightarrow\infty).
\]
The seminorm $p_1$ is continuous on Schwartz space.
If the support of $g$ lies in a translate of an open
$R$-cube, its distinct period translates are disjoint and
$|g_\eta|_\eta\le\|g\|_2$.

By \eqref{eq:slab-gaussian-domination},
\begin{equation}\label{eq:limit-cutoff-G}
 \E_{N,T}\exp(F^\sharp(g))
       \le\exp(|g_\eta|_\eta^2/2)\le\exp(p_R(g)^2/2)
                          \qquad(g\text{ real}).
\end{equation}

\begin{proposition}\label{prop:vtm-periodic-moments}
Suppose $N_i,T_i\to\infty$ and the laws $P_{N_i,T_i}$ of
$F^\sharp$ converge weakly on $H_{-4}$ to $\mu$.
Then $\E_\mu\exp(F(g))\le\exp(\|g\|_2^2/2)$ for every real
Schwartz two-form $g$. Every joint moment of a fixed finite
family of Schwartz tests converges along the same sequence.
\end{proposition}

\begin{proof}
Truncate $\exp(F(g))$ above, pass this bounded continuous
function to the limit, and remove the truncation.
Since $R_i\to\infty$, \eqref{eq:limit-torus-norm} proves the
Gaussian bound.

For fixed complex tests $g_1,\ldots,g_s$ put
\[
 Z_i(z)=\E_{N_i,T_i}\exp\left(\sum_a z_aF^\sharp(g_a)\right),
 \qquad
 Z(z)=\E_\mu\exp\left(\sum_a z_aF(g_a)\right).
\]
The evaluations are continuous on $H_{-4}$, and
\[
 \E_{N_i,T_i}
 \left|\exp\left(\sum_a z_aF^\sharp(g_a)\right)\right|^2
 \le\exp\left(2p_1\left(\operatorname{Re}\sum_a z_ag_a\right)^2\right).
\]
This is uniform in $i$ and on compact $z$-sets.
Uniform integrability and weak convergence give pointwise
convergence of $Z_i$ to $Z$. The exponential bounds make
these functions entire and locally uniformly bounded.
Cauchy's estimates give equicontinuity on compact sets;
pointwise convergence is therefore locally uniform.
Cauchy's formula passes every derivative at zero to the limit,
which proves joint moment convergence.
\end{proof}

We also record the varying-test estimate used below.
If $g_1,\ldots,g_J$ are real and their supports lie in one
open period cube, the Gaussian moment bounds and H\"older
give
\[
 \E\prod_{a=1}^J|F^\sharp(g_a)|
                     \le c_J\prod_{a=1}^J\|g_a\|_2
\]
for a constant $c_J$ depending only on $J$.
The same estimate holds for $\mu$.
Thus replacing each factor from a bounded family of tests
by another test, with $L^2$ difference at most $\delta$,
changes the expectation of a degree-$J$ monomial by at most
$Jc_JB^{J-1}\delta$, where $B$ bounds their $L^2$ norms.
This follows by telescoping the $J$ slots. For any fixed
finite family of compactly supported translates, the
support condition holds once $R_i$ is sufficiently large.

\subsection{Realizing the chosen parent by finite tori}

Let $\mu_i$ and $\mu$ be the prescribed laws in
\cref{thm:vtm-main}. Fix a metric $d_{\rm w}$ for weak
convergence of probability measures on the separable Hilbert
space $H_{-4}$. At spacing $\eta_i=L^{-N_i}$, take the
defining volume sequence for $\mu_i$. Choose an element
$T_i$ of that sequence so large that
\[
 d_{\rm w}(P_{N_i,T_i},\mu_i)<1/i,\qquad
 T_i\ge i+T^*(e)+1.
\]
Then
\[
 d_{\rm w}(P_{N_i,T_i},\mu)
 \le1/i+d_{\rm w}(\mu_i,\mu)\longrightarrow0.
\]
Both cutoffs diverge. This is the prescribed-parent
diagonal of \cite[Corollary~7.4]{ContinuumExistence},
with the finite-periodic moment conclusion supplied by
\cref{prop:vtm-periodic-moments}. Discard finitely many terms
so that $N_i\ge N_{\rm ref}(e)$.
At this fixed charge the translation spacing
$\lambda=R^*(e)/L$ is independent of $i$.

The four reference quantities
$c_{2,i},c_{4,i},u_{2,i},u_{4,i}$ are bounded by the estimates
of \cref{prop:transfer-UV-inputs}. A subsequence on which
they converge still has weak limit $\mu$.
For every fixed finite $Y\subset\lambda\Z^3$, the points
are eventually incongruent modulo the torus and all the
translated supports fit in a single open period cube.
The next section passes the Wick products along this
same coefficient subsequence. The threshold for support
containment may depend on $Y$; the coefficients and parent
law do not.

